% EEG-to-CLIP Alignment on THINGS-EEG — Paper Draft
% Target: NeurIPS ML4H Workshop or IEEE EMBC
% Author: Anurag Sharma, Galgotias University

\documentclass[conference]{IEEEtran}
\usepackage{amsmath,graphicx,booktabs,xcolor}
\usepackage{float}
\usepackage{placeins}
\usepackage{hyperref}
\usepackage{cite}

\begin{document}

\title{Contrastive EEG-to-CLIP Alignment on THINGS-EEG:\\
An Ablation of Target Modality, Temperature,\\
and Channel Count}

\author{
  \IEEEauthorblockN{Anurag Sharma}
  \IEEEauthorblockA{
    Department of Computer Science and Engineering
    (AI \& ML)\\
    Galgotias University, Greater Noida, India\\
    sharma.anurag0706@gmail.com
  }
}

\maketitle

% ──────────────────────────────────────────────
\begin{abstract}
We study the alignment of electroencephalography
(EEG) signals with pretrained CLIP embeddings
on the THINGS-EEG benchmark. Our contributions
are threefold. First, we present an
attention-based temporal-spatial encoder that
maps 17-channel EEG epochs to 512-dimensional
CLIP-compatible embeddings via bidirectional
InfoNCE contrastive loss. Second, we conduct
a systematic ablation over the contrastive
target modality (text CLIP vs.\ image CLIP
vs.\ a mixture), the training temperature
$\tau$, and the input channel count
(17 vs.\ 63). We find that image CLIP targets
outperform text targets, that mixing the two
modalities degrades performance below either
alone, and that expanding from 17 to 63
channels reduces accuracy under fixed
training-pair budget, implicating data scale
as the primary bottleneck. Our best
configuration attains a Top-5 retrieval
accuracy of 4.5\% on 200 held-out concepts
(1.8$\times$ the 2.5\% chance rate), a weak
but non-trivial alignment signal that we
characterise carefully with respect to
statistical significance. Third, we perform
a post-hoc channel-importance analysis of
the learned attention gates, finding that
the left parietal electrode P3 receives
the highest weight, consistent with the
P300 event-related potential and its
association with attentional salience.
We do not attempt image \emph{reconstruction}:
EEG embeddings are used for nearest-concept
\emph{retrieval} only, and image synthesis is
demonstrated by conditioning Stable Diffusion
on the retrieved concept's text label rather
than on EEG features directly. An extension
of this retrieval pipeline to sleep-EEG
settings is outlined as future work.
\end{abstract}

\begin{IEEEkeywords}
EEG, neural decoding, CLIP, contrastive
learning, brain-computer interface,
attention mechanism, retrieval
\end{IEEEkeywords}

% ──────────────────────────────────────────────
\section{Introduction}

Decoding visual content from human brain
activity has advanced rapidly in recent years,
with fMRI-based systems producing
high-fidelity reconstructions of viewed
images~\cite{takagi2023,mindvis2023}.
However, fMRI requires clinical
infrastructure, which precludes deployment
in naturalistic or longitudinal settings.
Electroencephalography (EEG) offers a
wearable alternative at the cost of
substantially lower spatial resolution and
signal-to-noise ratio.

Prior EEG decoding work has focused on the
\emph{awake} visual response: subjects view
images while EEG is recorded, and models are
trained to retrieve or reconstruct the viewed
image~\cite{scotti2023}. The dominant
technique is contrastive alignment to a
pretrained vision-language embedding, such
as CLIP~\cite{radford2021}, which provides
a semantically structured target space
without requiring a decoder. However, the
practical performance of these systems is
sensitive to several design choices ---
target modality, contrastive temperature,
input channel count, and available training
data --- that have not been systematically
ablated in the published literature.

\textbf{Contributions.}
We present an ablation study of these design
choices on the THINGS-EEG
benchmark~\cite{gifford2022}. Specifically:

\begin{enumerate}
  \item We describe an attention-based
  temporal-spatial EEG encoder that projects
  17-channel EEG epochs into CLIP-compatible
  512-dimensional embeddings.

  \item We ablate five training configurations
  spanning target modality (text CLIP, image
  CLIP, mixture), contrastive temperature
  ($\tau \in \{0.03, 0.05, 0.07\}$), and
  input channel count (17 vs.\ 63), reporting
  Top-1, Top-5, and Top-10 retrieval accuracy
  on 200 held-out concepts.

  \item We provide a statistical analysis of
  the best configuration against the chance
  baseline and identify data scale --- not
  architecture --- as the dominant limiting
  factor.

  \item We analyse the learned channel-attention
  gates and find that P3 (left parietal)
  carries the largest gate weight, consistent
  with the P300 ERP component and prior
  findings on attentional salience.

  \item We outline an extension of the retrieval
  pipeline to sleep-EEG settings as future
  work, explicitly noting that reconstruction
  of dream imagery is \emph{not} attempted in
  this paper.
\end{enumerate}

% ──────────────────────────────────────────────
\section{Related Work}

\subsection{EEG-to-Image Decoding}
Scotti et al.~\cite{scotti2023} demonstrated
EEG-to-image retrieval on THINGS-EEG using
contrastive alignment, achieving approximately
22\% Top-5 accuracy when training on
individual-subject EEG with GPU-scale batch
sizes. Their result establishes an upper
bound for the present study. Our work differs
in scope rather than in method: we do not
attempt to match this benchmark, and we
explicitly study how performance degrades
when training data is reduced through
subject-averaging, as required by
CPU-only training environments.

\subsection{Contrastive EEG Alignment}
CLIP~\cite{radford2021} demonstrated that
text and image representations can be aligned
in a shared embedding space via contrastive
learning on paired data. This formulation
has since been applied to a wide range of
non-visual modalities, including EEG. We
adopt the standard bidirectional InfoNCE
formulation and treat the CLIP image
embeddings as fixed targets, training only
the EEG encoder.

\subsection{Attention-based EEG Encoders}
Transformers~\cite{vaswani2017} and their
attention variants have become the default
architecture for sequence modelling in EEG
decoding. Our encoder applies multi-head
self-attention across channel tokens, a
temporal-spatial convolutional branch to
capture event-related dynamics, and a
gating mechanism that exposes per-channel
importance weights for post-hoc
interpretation. Related channel-attention
formulations have been explored for
sleep staging and motor imagery; we
adapt the pattern to the image-retrieval
task.

% ──────────────────────────────────────────────
\section{Methods}

\subsection{Dataset: THINGS-EEG}
We use the THINGS-EEG dataset~\cite{gifford2022},
comprising EEG recordings from 50 subjects
viewing 1{,}654 training and 200 test images
from the THINGS object concept
database~\cite{things2019}. Preprocessed
data provides 17-channel EEG epochs at
100\,Hz covering $-$200\,ms to 790\,ms
relative to image onset, with 80 repetitions
per concept per subject.

\textbf{Preprocessing.}
For each subject and each concept, we average
across the 80 repetitions, and then average
across the 50 subjects, yielding one clean
EEG epoch per concept: $(1654, 17, 100)$ for
training and $(200, 17, 100)$ for test.
Each channel is independently normalised to
zero mean and unit variance. Subject-averaging
is a computational necessity for our CPU-only
training environment; its effect on
performance is discussed in
Section~\ref{sec:disc}.

\subsection{EEG Encoder}

\textbf{Channel embedding.}
Each of the 17 EEG channels (100 timepoints)
is projected independently via a shared
linear layer to a $d_\text{model}{=}128$-dimensional
token, producing a sequence of channel
tokens with shape $(B, 17, 128)$.

\textbf{Channel-wise attention.}
Multi-head self-attention (4 heads) is
applied across the 17 channel tokens. A
sigmoid gating mechanism produces
per-channel importance weights, enabling
post-hoc interpretation of electrode
contributions.

\textbf{Temporal-spatial convolution.}
Two parallel convolutional paths capture
complementary structure: a temporal path
(Conv1D across time within each channel,
kernel size 15) capturing ERP dynamics,
and a spatial path (Conv1D across channels
at each timepoint, kernel size 5) capturing
electrode co-activation patterns. The two
paths are fused via a linear layer.

\textbf{Aggregation and projection.}
The channel-attended representation and
the temporal-spatial representation are
concatenated and projected to a
512-dimensional CLIP-compatible embedding,
which is L2-normalised. The encoder contains
626{,}449 trainable parameters.

\subsection{CLIP Target Computation}

For each training concept, we compute a CLIP
image embedding~\cite{radford2021} by loading
all available images from the corresponding
THINGS concept folder, computing the CLIP
ViT-B/32 image feature for each, and
averaging across images before
re-normalisation. This yields a robust
prototype embedding per concept. The mean
inter-concept cosine similarity is 0.529
for image targets, compared with 0.713 for
text-based targets of the form
\texttt{``a photo of a \{concept\}''}, confirming
that image targets provide a more
discriminative learning signal.

\subsection{Training}

\textbf{Loss function.}
We use bidirectional InfoNCE contrastive loss
with temperature $\tau$:
\begin{equation}
\mathcal{L} = -\frac{1}{2N}\sum_{i=1}^{N}
\left[
  \log\frac{e^{\mathbf{z}_i \cdot
    \mathbf{c}_i/\tau}}
    {\sum_j e^{\mathbf{z}_i \cdot
    \mathbf{c}_j/\tau}}
  + \log\frac{e^{\mathbf{c}_i \cdot
    \mathbf{z}_i/\tau}}
    {\sum_j e^{\mathbf{c}_i \cdot
    \mathbf{z}_j/\tau}}
\right]
\end{equation}
where $\mathbf{z}_i$ is the normalised EEG
embedding and $\mathbf{c}_i$ is the
normalised CLIP image embedding for concept
$i$ within the mini-batch.

\textbf{Optimiser and schedule.}
AdamW with weight decay 0.05, OneCycleLR
schedule (peak learning rate
$5 \times 10^{-4}$, 10\% warmup),
200 epochs, batch size 128.
EEG augmentation during training: additive
Gaussian noise ($\sigma{=}0.08$), random
channel dropout (2--3 channels, $p{=}0.25$),
and temporal shift ($\pm 3$ timepoints,
$p{=}0.30$).

\subsection{Retrieval Protocol}
At evaluation, each of the 200 test EEG
epochs is encoded to a 512-dimensional
embedding, and cosine similarity is computed
against the 200 test-concept CLIP embeddings.
We report Top-1, Top-5, and Top-10 accuracy,
with the chance rates $1/200$, $5/200$, and
$10/200$ respectively. All reported accuracies
are on the held-out test concepts.

Fig.~\ref{fig:pipeline} illustrates the
overall pipeline.

% ── FLOAT 1: FIGURE 1 ──
\begin{figure*}[!t]
\centering
\includegraphics[width=\textwidth]{phase2/fig1_pipeline}
\caption{Overview of the EEG-to-CLIP retrieval
pipeline. EEG epochs are encoded by the
proposed attention-based temporal-spatial
encoder into a 512-dimensional CLIP-compatible
embedding, and retrieved by cosine similarity
against the CLIP embeddings of the 200 test
concepts. Image synthesis, shown for
illustration only, is conditioned on the
retrieved concept's \emph{text} label, not
on the EEG embedding directly.}
\label{fig:pipeline}
\end{figure*}
\FloatBarrier
% ── END FLOAT 1 ──

% ──────────────────────────────────────────────
\section{Experiments}

\subsection{Ablation: Target Modality and Temperature}

Table~\ref{tab:ablation} reports Top-5
retrieval accuracy on the 200 test concepts
under five experimental conditions. All
configurations use the same encoder, optimiser,
and training schedule; only the target
modality, temperature, or channel count
varies.

% ── FLOAT 2: TABLE 1 ──
\begin{table}[H]
\centering
\caption{Ablation over target modality,
contrastive temperature, and channel count.
Test set: 200 THINGS-EEG concepts. Chance
Top-5 = 0.025. Best configuration in bold.}
\label{tab:ablation}
\begin{tabular}{lcc}
\toprule
\textbf{Configuration} &
\textbf{Top-5} &
\textbf{$\times$Chance} \\
\midrule
Chance & 0.025 & 1.0$\times$ \\
\midrule
Text CLIP, $\tau{=}0.07$ & 0.035 & 1.4$\times$ \\
Mixed (70\% img), $\tau{=}0.05$ & 0.030 & 1.2$\times$ \\
Image CLIP, $\tau{=}0.03$ & 0.040 & 1.6$\times$ \\
\textbf{Image CLIP, $\tau{=}0.05$} & \textbf{0.045} & \textbf{1.8$\times$} \\
Image CLIP, 63-ch, $\tau{=}0.05$ & 0.035 & 1.4$\times$ \\
\bottomrule
\end{tabular}
\end{table}
% ── END FLOAT 2 ──

Image targets outperform text targets by
0.010 Top-5 accuracy (28.6\% relative),
consistent with their higher inter-concept
discriminability (mean cosine similarity
0.529 vs.\ 0.713). Mixing the two target
modalities \emph{reduces} performance below
the text-only baseline, suggesting that text
and image CLIP representations occupy
sufficiently distinct regions of embedding
space that their mixture introduces noise
rather than combining signal.

Increasing the channel count from 17 to 63
reduces performance under our fixed
training-pair budget. The expanded input
dimensionality (6{,}300 vs.\ 1{,}700 features
per epoch) cannot be constrained by 1{,}654
training pairs, leading to overfitting and
training instability. This identifies data
scale, rather than encoder capacity, as the
primary bottleneck.

Fig.~\ref{fig:ablation} visualises Top-5
accuracy across all configurations.

% ── FLOAT 3: FIGURE 2 ──
\begin{figure}[H]
\centering
\includegraphics[width=\columnwidth]{phase2/fig2_ablation.png}
\caption{Top-5 retrieval accuracy across the
five configurations in Table~\ref{tab:ablation}.
Image CLIP targets with $\tau{=}0.05$ yield the
best result. Mixing text and image targets
degrades performance relative to either
modality alone.}
\label{fig:ablation}
\end{figure}
% ── END FLOAT 3 ──

\subsection{Channel Importance Analysis}

The learned attention gates expose which
electrodes contribute most to the final
embedding. Under the 17-channel configuration,
P3 (left parietal) receives the highest gate
weight. P3 is the scalp projection of the
P300 ERP component, which is associated with
attentional orientation toward salient
stimuli. Under the 63-channel configuration,
temporal-region electrodes (T7, T8, and
neighbours) dominate, consistent with the
role of inferior temporal cortex in visual
object identity processing --- the terminal
stage of the ventral visual
stream~\cite{dicarlo2012}.

Fig.~\ref{fig:attention} shows the learned
channel importance weights.

% ── FLOAT 4: FIGURE 3 ──
\begin{figure}[H]
\centering
\includegraphics[width=\columnwidth]{phase2/fig3_channel_attention}
\caption{Learned channel gate weights,
averaged across the 200 test EEG epochs.
P3 (left parietal) receives the highest
weight, consistent with the P300 ERP
component's association with attentional
orientation toward visual stimuli.}
\label{fig:attention}
\end{figure}
% ── END FLOAT 4 ──

\subsection{Qualitative Analysis}

Retrieval errors exhibit a consistent
pattern: the retrieved concepts tend to
share low-level visual structure with the
true concept rather than its semantic
category. The clearest example in our test
set is the retrieval of \textit{blowtorch}
at rank 3 for the true concept \textit{baton};
both are elongated cylindrical objects with
similar aspect ratios and silhouettes. We
present this as a single illustrative
observation rather than as evidence for a
general shape-encoding hypothesis.

Fig.~\ref{fig:qualitative} shows
representative generation outputs.

% ── FLOAT 5: FIGURE 4 ──
\begin{figure}[H]
\centering
\includegraphics[width=\columnwidth]{phase2/fig4_generated_examples}
\caption{Illustrative synthesis outputs.
Images generated by Stable Diffusion
v1.5~\cite{rombach2022} conditioned on the
\emph{text labels} of the top-1 retrieved
concepts (french horn, lamp, boat). These
are shown to illustrate the downstream
application; the EEG embedding conditions
the retrieval, not the synthesis.}
\label{fig:qualitative}
\end{figure}
% ── END FLOAT 5 ──

\subsection{Statistical Significance}
\label{sec:stats}

The best configuration attains 4.5\% Top-5
accuracy on 200 test concepts, corresponding
to 9 correct retrievals out of 200 against
a chance rate of 2.5\% (expected 5 correct).
A one-sided exact binomial test against
the chance rate yields $p \approx 0.05$,
which places the result at the boundary of
conventional statistical significance. We
therefore characterise our result as a
\emph{weak but non-trivial} alignment signal
rather than as competitive retrieval
performance. This framing is important:
with only 1{,}654 training pairs, a
contrastive objective has limited capacity
to shape a well-separated 512-dimensional
embedding space, and the observed gap
between our result and the 22\% Top-5
reported by Scotti et al.~\cite{scotti2023}
is consistent with their use of
individual-subject training (roughly $10\times$
more pairs) and GPU-scale batch sizes
(larger effective negative sets per step).

% ──────────────────────────────────────────────
\section{Discussion}
\label{sec:disc}

\subsection{Comparison to Prior Work}

Scotti et al.~\cite{scotti2023} achieved
approximately 22\% Top-5 accuracy on
THINGS-EEG using individual-subject EEG
(rather than subject-averaged) with
GPU-scale batch training. Our 4.5\% Top-5
reflects two primary differences: (1) we
average EEG across subjects, reducing the
effective number of training pairs from
roughly 16{,}540 to 1{,}654; and (2) our
CPU-only environment limits the batch size
to 128, reducing the number of negative
pairs per update from which contrastive
learning derives its signal. With
individual-subject training and GPU access,
we anticipate performance approaching
published benchmarks; we regard our
contribution as a controlled ablation of
the design choices that mediate this gap,
not as a new state-of-the-art.

\subsection{What This Paper Does and Does Not Show}

\textbf{Does:} We show that on the
THINGS-EEG benchmark with subject-averaged
EEG, image CLIP targets provide a
measurably better contrastive signal than
text CLIP targets; that mixing the two
modalities degrades performance; that a
63-channel input is not beneficial under
a fixed training-pair budget; and that the
learned channel gates concentrate on P3,
a physiologically interpretable electrode.

\textbf{Does not:} We do not reconstruct
images from EEG. We do not train on sleep
EEG. We do not claim state-of-the-art
retrieval. The images shown in
Fig.~\ref{fig:qualitative} are generated
from the retrieved concept's \emph{text
label}, not from the EEG embedding directly;
the EEG signal contributes to the pipeline
only via the retrieval step. No closed-loop
brain-computer interface has been built or
tested.

\subsection{Limitations}

Three limitations bound the scope of our
conclusions. First, our training data are
EEG recorded during \emph{awake} visual
stimulation, while the eventual target
application is sleep-EEG decoding; the
domain shift between these regimes has not
been quantified and may be substantial.
Second, subject-averaging reduces the
effective training-pair count by an order
of magnitude relative to published
individual-subject results, and this is
the dominant reason our accuracy is
modest. Third, our qualitative results
rest on a small number of illustrative
examples rather than on a systematic
error analysis.

\subsection{Future Work}

We plan to extend this work in three
directions. First, we will train on
individual-subject EEG to recover the
$\sim 10\times$ training-pair advantage
and approach published benchmarks.
Second, we will integrate a consumer
EEG headband (e.g., Muse S) to enable
real-time processing and to test the
pipeline on non-laboratory data. Third,
and most speculatively, we will investigate
whether the alignment learned on awake
EEG transfers to sleep-EEG settings,
with the long-term goal of conditioning
image synthesis on internally generated
imagery during REM sleep. We emphasise
that this last direction is a research
proposal, not a demonstrated capability.

% ──────────────────────────────────────────────
\section{Conclusion}

We have presented a controlled ablation of
EEG-to-CLIP contrastive alignment on the
THINGS-EEG benchmark with subject-averaged
training data. Image CLIP targets
outperform text targets, mixing target
modalities is harmful, and expanding input
channel count is not beneficial at the
1{,}654-pair training scale. Our best
configuration attains 4.5\% Top-5
retrieval accuracy (1.8$\times$ chance),
at the boundary of statistical significance,
and we identify data scale as the primary
bottleneck. Channel-attention analysis
reveals physiologically interpretable
importance concentrated on P3. Extensions
to individual-subject training, consumer
EEG hardware, and sleep-EEG decoding are
outlined as future work; no reconstruction
from sleep EEG is attempted here.


% ──────────────────────────────────────────────
\bibliographystyle{IEEEtran}
\bibliography{references}
\end{document}